\begin{tablalafrox}{Dotación, estabilización y soporte entre etapas. Fuente: elaboración a partir de Distribuidora Puelche S.A. (2026a, art. 17.3, p. 13); Distribuidora Puelche S.A. (2026b, sec. 19.2 y cap. 18, pp. 33–35); LafroX (2026d, T-14); LafroX (2026e, T-15).}{tab_T18_prevision}
  {L{3.5cm} L{3.5cm} Y}
  {\cab{Período} & \cab{Previsión de puestos simultáneos E1} & \cab{Tratamiento}}
  Meses 13–15 & 13 & Acompañamiento de marcha blanca, sin reducción que comprometa volumen o evidencia \\
  Primeras cuatro semanas posteriores a H7, mes 16 & 13 & Estabilización reforzada; no se absorbe por desarrollo E2 \\
  Mes 17 & 7: dos CD, tres plataformas, una ruta y coordinación & Sólo si todos los usuarios y transacciones del alcance están certificados/registrados, conciliación sin diferencias inexplicadas y cero incidentes críticos/altos \\
  Meses 18–20 & 3: dos CD y coordinación; especialistas de ruta/plataformas bajo demanda & Sólo con criterios anteriores sostenidos; conserva NOC/mesa y reserva DES/QA por separado \\
  Primeras cuatro semanas tras H12, meses 21–22 & 13 puestos de acompañamiento E2 como supuesto conservador & Con el acta calculada el 25 de octubre de 2028, el acompañamiento va del 26 de octubre al 22 de noviembre; bodega y ruta verifican no degradación de E1. Se recalibra con el alcance efectivo antes de H11 y con la fecha efectiva del acta \\
\end{tablalafrox}
